\section{Statement, dependence, and antecedents}\label{sec:scope}
Let $\Sigma$ be a finite alphabet with an involution $\cm:\Sigma\to\Sigma$.
For a word $s=s_1\cdots s_h$, put
$\rho(s)=\cm(s_h)\cdots\cm(s_1)$ and $\rho(\eps)=\eps$.
A word is feasible for a finite listed family $\mathcal S$ if, for every
$s\in\mathcal S$, it contains $s$ or $\rho(s)$ as a contiguous substring.
Write $\OPT(\mathcal S)$ for its minimum symbol length. Alphabet labels,
the involution table, word symbols, and list boundaries are explicit;
$N$ denotes their total encoded bit length.

\begin{theorem}[Source-relative approximation]\label{thm:main}
Assume the structural results of \cite{ordinary} at revision
\textnormal{\texttt{fd4aeeb2ee4fc729c18d98444fed42fd0529eeeb}}, with the contracts
specified below. There is a deterministic algorithm, polynomial in $N$,
which for every such $(\Sigma,\cm,\mathcal S)$ returns a feasible word $T$
with $|T|\le2\OPT(\mathcal S)$. The assertion includes arbitrary fixed
letters and words, empty words and the empty list, duplicates, and substring
redundancy. Neither the alphabet nor the word lengths are bounded constants.
\end{theorem}

The source's truths are assumptions of this article. We do not audit them
or derive an independent proof of the ordinary superstring theorem.
In particular, an unspecified ordinary factor-two algorithm would not
supply the contracts used here. Yamano--Shibuya's reduction in
\cite[Theorem 2]{yswabi} also explains why an unrestricted reverse-complement
result cannot be advertised as bypassing the ordinary problem.

\subsection{The precise source interface}\label{sec:contracts}
The following labels identify results in the pinned source, not labels
of this article. Its section files are under \texttt{build/sections/}.
\begin{enumerate}
\item \path{02-counts.tex}: the blocking functional, extension and
ordinary period rules, and the occurrence and balance proofs
\path{lem:counts}, \path{lem:base}. Their use requires a nonnegative
count vector on a substring-closed set satisfying those rules. It does
not require choosing the smallest ordinary count vector.
\item \path{03-layers.tex}: \path{lem:recorded}, \path{lem:layers},
\path{lem:matching-window}, \path{lem:agreement}, and
\path{lem:max-rotation}. These provide recorded counts, sorted primitive
passages, actual matching windows, and the two periodic comparison bounds.
\item \path{04-connections.tex}: \path{lem:ordered-windows} and
\path{lem:band}; the \emph{proof} of \path{lem:collective}; and
\path{lem:host} after its hypotheses are established on a single host
path. We state the relevant contracts in Section~\ref{sec:operations}.
The source collective statement has $k$ layers and $K=kp$; our odd-core
use has $2k-1$ layers and $K=kp$, and uses the proof's geometric inequalities.
\item \path{05-processing.tex}: its ordinary threshold, partition,
and individual-link/request construction, applied to nonself text orbits.
\item \path{06-cycles.tex}: the proof of \path{lem:cycle-contact}
and the free-period strategy of \path{lem:free-period} and
\path{thm:connection}. We adapt the geometric proof to odd cores and
justify the two-port quotient separately; the original one-sided
connection theorem is not assumed equivariant.
\item \path{07-algorithm.tex}: finite record search, explicit coordinate
representation, and the suffix invariant for spelling the output.
\end{enumerate}
All these hypotheses and their symmetric applications are made explicit
below. Source-paid arguments are recalled in our notation where needed to
make the adaptation checkable.

\subsection{Related operations and the credit boundary}\label{sec:antecedents}
The arbitrary-involution formulation, formal copies of fixed words,
paired cycle selection, both-orientation periodic comparisons, and an
$8/3$ approximation occur in Yamano--Shibuya \cite{yswabi}, especially
Section 2, Lemmas 7, 9, 12, and Theorems 11--13. Their CPM paper
\cite{yscpm} proves $3.75$ and $2.875$ guarantees and, in Lemma 18,
compares an optimum with its entire reverse complement. Thus the generic
use of $T$ and $\rho(T)$ is an antecedent, not a contribution here.
Fici et al.\ \cite{fici} already use distinct palindrome vertex copies
and paired reverse paths in Sections 2--3, and reversal surgery in
Lemma 6. Their Theorem 7 is a one-half \emph{compression} guarantee;
it is not a factor-two length guarantee.

Jiang, Li and Du studied superstrings with flipping, as described by
Fici et al.\ and Yamano--Shibuya \cite{jld,fici,yscpm,yswabi}.
The 1992 full text was unavailable in our bounded searches, so our
comparison uses those later accounts. Fici et al.\ report a $4$ length
bound for a greedy-like reversal algorithm. The CPM 2026 account
distinguishes the $4$-approximate MGREEDY extension from a suggested
$3$-approximate TGREEDY extension; WABI 2026 summarizes the earlier bound
as $3$. We do not assign an independent theorem statement to the unread
1992 body.

Goldberg--Karzanov's symmetric decomposition and parity results
\cite[Theorem 3.1 and Corollary 3.2]{gk} precede the pairing operations
below. Their Section 5 already exchanges complementary halves of two
self-symmetric cycles meeting at two mate vertices. We credit that splice;
our shared-axis lemma supplies the meeting hypothesis for possibly
different ordered staircase profiles. Their setting has no fixed vertices,
whereas our hierarchical graph has fixed vertices, including $\eps$;
we therefore give the needed parity and traversal proof explicitly.
Bidirected incidence walks and balanced connected Euler criteria also
appear in Medvedev et al.\ \cite[Section 2.2, Lemma 1]{medvedev}
and Kundeti et al.\ \cite[Section 6]{kundeti}. Their spectrum-constrained
cyclic assembly results are not imported as exact algorithms for our
arbitrary-length linear requirements.

Extremal and critical rotations are established tools in
Breslauer--Jiang--Jiang \cite{bjjjournal,bjjreport} and Mucha \cite{mucha}.
We cite the 1997 journal article as the archival reference, but our
operational comparison uses the 1996 BRICS report RS-96-21; the journal
body was not read, and report lemma numbers are not transferred to it.
The report's concluding remarks credit Crochemore with a Lyndon
minimum/maximum-rotation argument. The hierarchical representation and
the separation between cycle-cover bounds and connectivity also precede
this work; see the Collapsing framework \cite[Sections 2.2 and 7.2]{collapsing}.

The contribution claimed here is exactly one matched package:
the phase-free self-text count rule in Lemma~\ref{lem:domination}, the transfer
to adjacent sorted profiles sharing both axes in Lemma~\ref{lem:axes}, and
the $2k-1$ odd core inside a $2k$-layer unit in Lemma~\ref{lem:odd-core}.
No separate novelty is claimed for phase parity, tagged counts,
maxima or even rounding, decomposition, splicing, anchors, batch flips,
request reclassification, rotations, cycle accounting, half-Euler
extraction, or complexity bookkeeping.

\subsection{Why the matched package is needed}\label{sec:motivation}
For identity complement on $\{b,c\}$, take
$\mathcal S=\{bbcbbc,bcbbcb,cbbcbb\}$. The ordinary rules symmetrized
by orbit maxima and even rounding give one-letter counts $4,2$, hence
$W_{\rm sym}=6$. But $\OPT=7$: a length-six word cannot contain the two
distinct length-six orientation orbits, and $bbcbbcb$ satisfies both.
This refutes only a connection theorem that would retain those counts
and yield output length at most the same $W_{\rm sym}$. It does not
refute factor two or all equivariant methods. The stronger rule below
gives $W=12$ on this instance.

A different diagnostic, $\{aaa,aab,abcaa,accc\}$ with $a\leftrightarrow b$
and $c$ fixed, has optimum $10$. Minimum constrained-cycle-cover
orientation ties can lead to ordinary forced-count certificates $11$
or $9$. This shows that arbitrary minimum tie-breaking does not ensure
the desired $W$ certificate; it does not exclude all orientation-selection
methods. These finite certificates motivate the construction and are not
used to prove Theorem~\ref{thm:main}.

Boosting only fixed words is also insufficient for the intended proof:
in $(abcacb)^{\ZZ}$, the layer $z(x)=x+8$ has fixed-axis windows at starts
$0,3$, whereas its unique greatest forward rotation starts at $4$.
The two windows there are $cbabcac$ and $cbabcacb$, neither fixed.
A count rule confined to fixed-axis starts would not justify the
maximal-rotation cycle argument. The rule below holds at every phase
of a self-reversed text.
